% TOP_PROBLEM: {"id": 1600012, "title": "Grothendieck–Katz p-curvature conjecture", "category_id": 4, "category": {"id": 4, "name": "algebra", "display_name": "Algebra", "description": "Group theory, ring theory, field theory, and algebraic structures.", "slug": "algebra", "order_index": 4, "created_at": "2026-07-31T15:26:25.671Z"}, "status": "open"}
% FIRST_POSTED: null
% ENTRY_KIND: statement_only
% Imported from: batch06.tex; UnsolvedMath ID 1600012
% Compile twice: lualatex 1600012.tex
\documentclass[11pt]{article}
\usepackage{fontspec}
\setmainfont{Latin Modern Roman}
\usepackage{amsmath,amssymb,mathtools,unicode-math}
\setmathfont{Latin Modern Math}
\usepackage{xurl,hyperref}
\hypersetup{colorlinks=true,urlcolor=blue,pdftitle={UnsolvedMath Problem 1600012}}
\newcommand{\sref}[2]{\hyperlink{source:#1:#2}{[S#2]}}
\newcommand{\cataloguescope}[1]{\paragraph{Mathematical question.}#1}
\begin{document}
% UnsolvedMath ID 1600012; source catalogue code AMR-015-0012
\setcounter{section}{1600011}
\section{Grothendieck–Katz p-curvature conjecture}\label{1600012}
{\small\sffamily UnsolvedMath ID 1600012 · Algebra}
\subsection{Definitions and mathematical statement}

\paragraph{Shared notation.}$\mathbb N=\{1,2,\ldots\}$, and $\mathbb Z,\mathbb Q,\mathbb R,\mathbb C$ denote the integers, rationals, reals and complex numbers. Any other conventions are specified locally.


Let $X$ be a smooth connected complex algebraic variety and $(V,\nabla)$ an algebraic vector bundle with integrable connection. Choose a model of $X$ and $(V,\nabla)$ over a finitely generated $\mathbb Z$-subalgebra $R\subset\mathbb C$, smooth after localizing $R$. At a good reduction $\mathfrak p$ of residue characteristic $p>0$, define its $p$-curvature by
\[\psi_{\mathfrak p}(D)=\nabla_D^p-\nabla_{D^{[p]}},
\]
where $D$ is a derivation of the reduced structure sheaf and $D^{[p]}$ is its $p$-fold iterate. Vanishing means this endomorphism is zero for every $D$. A full set of algebraic horizontal solutions means that some finite étale cover $Y\to X$ makes the pulled-back connection trivial.
\paragraph{Statement recorded in the supplied source.}
If the $p$-curvature vanishes for almost all good reductions of the model, must $(V,\nabla)$ have a full set of algebraic horizontal solutions? Here almost all means after excluding finitely many residue characteristics (and localizing the arithmetic model). \sref{1600012}{2}
\subsection{Short English statement}
Do solutions modulo almost every prime force the original differential equation to have a full set of algebraic solutions?
\subsection{Sources}
\begin{description}
\item[\hypertarget{source:1600012:1}{S1}] ULAM.AI and the UnsolvedMath Contributors, UnsolvedMath: A Curated Collection of Open Mathematics Problems, record 1600012 (AMR-015-0012). CC BY 4.0. \url{https://huggingface.co/datasets/ulamai/UnsolvedMath}.
\item[\hypertarget{source:1600012:2}{S2}] B. Farb and M. Kisin, \emph{The Grothendieck--Katz Conjecture for certain locally symmetric varieties}, arXiv:0807.1152 (2008), §1, Conjecture 1 and the preceding arithmetic model. \url{https://arxiv.org/abs/0807.1152}.
\end{description}
\label{end:1600012}
\end{document}
